% ================== Chapitre 7 : hypotheses, limites, conclusion ============
\chapter{Hypothèses, limites et points à valider}
\label{ch:limites}

La qualité d'un dossier de conception tient autant à ce qu'il établit qu'à ce
qu'il signale comme non établi. Les éléments ci-dessous distinguent les
hypothèses prises faute de donnée, les points dont la vérification conditionne
l'application, et les limites structurelles que la conception ne peut pas lever.

\section{Hypothèses formulées}

\begin{longtable}{@{}R{0.27}R{0.58}R{0.15}@{}}
\caption{Hypothèses et choix à confirmer}\label{tab:hypotheses}\\
\hautfilet
\textbf{Élément} & \textbf{Justification} & \textbf{Statut}\\
\medfilet
\endfirsthead
\hautfilet
\textbf{Élément} & \textbf{Justification} & \textbf{Statut}\\
\medfilet
\endhead
\basfilet
\endfoot
Durée de bail DHCP $B = 3\,600$ s &
  Donnée manquante, relevant de T004. Elle fixe la durée de vie des
  enregistrements dynamiques et la quarantaine par l'inégalité
  \eqref{eq:ttl}, qui reste valable si la valeur change. &
  Hypothèse\\
Quarantaine $Q = 900$ s avant réutilisation d'une adresse &
  Le document 03 laisse cette durée \statut{à valider}. La valeur proposée
  satisfait \eqref{eq:ttl} avec une marge $\delta = 300$ s. &
  Proposition\\
Récursion complète vers la racine publique &
  Préférée à une redirection vers un résolveur d'opérateur, dont le
  comportement n'est pas maîtrisé et dont la disponibilité dépendrait d'un
  équipement non manageable. & Choix\\
VNI 30404 pour le VNet Services &
  Pris dans la réserve de travail 30404 à 30408 du document 02. L'identifiant
  réellement attribué par le contrôleur doit être relevé et rapproché de cette
  réserve. & Proposition\\
Moteur de stockage SQLite &
  Suffisant pour $N_Z = 113$ zones et un régime courant $N_R \approx 1\,330$
  enregistrements. Un moteur relationnel partagé reste une extension. &
  Choix\\
\end{longtable}

\section{Points à vérifier avant application}

\begin{itemize}
  \item \textbf{Versions exactes des logiciels.} Les directives \opt{primary} et
    \opt{secondary} supposent une version 4.5 ou supérieure du serveur
    autoritatif. La branche 5 du résolveur propose un format de configuration
    YAML qu'il faut distinguer du format classique employé ici. La disponibilité
    de la journalisation dnstap dépend également des options de compilation du
    paquet retenu.
  \item \textbf{Collecteur de journalisation.} Le composant qui consomme le flux
    dnstap et produit les fichiers lus par Monitoring n'est pas arrêté. Son
    choix doit respecter la règle d'absence d'initiation vers Monitoring.
  \item \textbf{Comportement réel de la sortie externe.} La récursion vers la
    racine suppose que le routeur R1 laisse passer le port 53 sortant en UDP et
    en TCP sans altération. Le document 01 rappelle que les fonctions réelles de
    ce routeur restent à vérifier.
  \item \textbf{Correspondance entre session VPN et tenant.} Tant qu'elle n'est
    pas fournie par le lot VPN et IAM, le pool VPN utilisateur reste privé de
    résolution interne, comme expliqué à la section~\ref{sec:vpn}.
  \item \textbf{Coût du script de politique.} Le script de la
    section~\ref{sec:script} s'exécute à chaque requête. Son incidence sur la
    latence doit être mesurée à la charge attendue avant mise en service, même
    si son contenu est volontairement minimal.
\end{itemize}

\section{Limites structurelles assumées}

\begin{itemize}
  \item La redondance du service couvre la perte d'une machine virtuelle ou
    d'un hôte. Elle ne couvre ni la perte du commutateur SW1, ni celle du
    routeur R1, qui restent des dépendances communes imposées par le cadrage
    matériel.
  \item La politique de vue repose sur l'adresse source. Elle suppose donc que
    l'usurpation d'adresse à l'intérieur d'un VNet est traitée par la
    segmentation et les protections du lot VNet et SDN. Le DNS ne constitue pas,
    à lui seul, un mécanisme d'authentification du demandeur.
  \item Un tenant peut énumérer ses propres noms. Cette propriété est inhérente
    à un service de résolution ouvert à ses clients et ne franchit aucune
    frontière de tenant.
  \item Aucun déploiement n'a été réalisé et aucune mesure n'a été produite. Les
    volumes, les délais et les capacités énoncés sont des estimations de
    conception, à confronter à la maquette avant toute mise en service.
\end{itemize}

\section*{Conclusion}
\addcontentsline{toc}{section}{Conclusion}

Ce dossier fournit ce que la tâche T005 demandait : une architecture de
résolution arrêtée, un espace de nommage entièrement dérivé du plan d'adressage
du document 02, un mécanisme démontré d'isolation entre domaines, des
configurations applicables pour les quatre instances, et un cycle de vie des
enregistrements compatible avec la machine d'états du document 03.

Trois décisions en constituent l'ossature. La séparation stricte des rôles
autoritatif et récursif, qui rend l'isolation applicable plutôt que seulement
souhaitée. L'indexation de la politique de vue sur l'étiquette $\tau(a)$
calculée avant le cache, qui empêche le cache de contourner la politique. Enfin
le propriétaire d'écriture unique, qui rend la publication rejouable et la
réconciliation possible.

Le service reste conditionné aux hypothèses de ce chapitre et à la simulation
préalable exigée par \exig{net-fr-020}. L'exécution des contrôles du
chapitre~\ref{ch:recette}, la mesure des cibles de reconstruction et l'archivage
des preuves relèvent de la tâche T009, qui interviendra lorsque les lots DHCP,
temps et agent seront disponibles.
